\subsection{Analyse préliminaire des risques}
\label{subsec:analyse_prelem}

L’analyse préliminaire des risques constitue le point de départ de la sûreté de fonctionnement. Elle consiste à examiner, dès
les premières phases de conception, les situations susceptibles de mettre en cause la sécurité du vol, la protection des
personnes et la conformité du lanceur aux exigences du cahier des charges Fusex bi-étage. Cette étape préliminaire permet
d’identifier les domaines techniques où un niveau de vigilance particulier est nécessaire et oriente les analyses détaillées
réalisées dans les sections suivantes. \\

Le projet SP-02 présente plusieurs caractéristiques critiques — séparation passive par aérofreinage, allumage conditionnel du
second étage en vol, structure en peau porteuse, mécanismes de récupération et avionique distribuée — qui introduisent des
risques spécifiques. L’analyse vise à repérer les situations potentiellement dangereuses associées à ces sous-systèmes et à
déterminer comment elles pourraient affecter l’intégrité et le pilotage du lanceur. \\

Cette première étape se conclut par l’identification des Événements Redoutés Principaux (ERP), c’est-à-dire des situations
finales dont la survenue compromettrait directement la sécurité ou la trajectoire du lanceur. Ils constituent la base de
structuration des sections~2.2 et~2.3.

\subsubsection{Événements redoutés principaux (ERP)}

Les Événements Redoutés Principaux (ERP) correspondent aux situations dangereuses qu’il est nécessaire d’éviter pour garantir
un vol sûr et conforme. Ils constituent les \textbf{états finaux indésirables} que des chaînes de défaillances pourraient
produire, et non les pannes ou dysfonctionnements eux-mêmes.\\

Leur rôle est de fournir une structure claire à l’analyse de sûreté : chaque mode de défaillance identifié ultérieurement sera
évalué en fonction de sa capacité à conduire à l’un de ces événements, et la criticité associée sera examinée dans l’AMDEC. Les
ERP définis pour le projet SP-02, en conformité avec la règle ASF1, sont présentés dans le tableau~\ref{tab:ERP_SP02}.

\begin{table}[h!]
\centering
\begin{tabular}{|c|p{12.5cm}|}
\hline
\textbf{ERP1} & 
\textbf{Déviation de trajectoire / sortie du gabarit.}\newline
Perte de contrôle ou trajectoire non conforme, provoquée par une instabilité,
une perturbation aérodynamique, un défaut de séparation, une rupture structurelle
ou tout événement compromettant la dynamique nominale de vol. 
\\ \hline

\textbf{ERP2} & 
\textbf{a - Allumage intempestif ou non autorisé du second étage en phase sol ou rampe.}\newline
Mise à feu du moteur du deuxième étage alors que la fusée n’est pas encore en vol 
ou dans un état non prévu pour supporter une propulsion, représentant un risque
majeur pour les personnes et les biens. \newline\newline
\textbf{b - Allumage intempestif ou non autorisé du second étage en vol.}\newline
Allumage en dehors des conditions requises (séparation non confirmée, attitude non acceptable,
fenêtre temporelle inactive, logique d’ordre erronée), pouvant entraîner une perte
de contrôle ou une violation du gabarit de vol.
\\ \hline

\textbf{ERP3} & 
\textbf{Non-déploiement ou déploiement dangereux du parachute (E1 ou E2).}\newline
Absence d’ouverture, ouverture tardive, partielle ou perturbée, conduisant
à une chute à vitesse excessive, susceptible de générer un impact dangereux.
\\ \hline

\textbf{ERP4} & 
\textbf{Rupture structurelle majeure entraînant une perte de contrôle.}\newline
Dégradation critique de la structure (aileron, peau porteuse, jonctions, coiffe)
pouvant provoquer une instabilité brutale, une rotation excessive ou 
une dispersion hors gabarit.
\\ \hline
\end{tabular}
\caption{Liste des Événements Redoutés Principaux (ERP) du projet SP-02}
\label{tab:ERP_SP02}
\end{table}